\usepackage[T1]{fontenc}
\usepackage[utf8]{inputenc}
\usepackage{lmodern}
\usepackage{amsmath,amssymb,amsthm}
\usepackage{bm}
\usepackage{booktabs}
\usepackage{multirow}
\usepackage{adjustbox}
\usepackage{graphicx}
\usepackage{xcolor}
\usepackage{tikz}
\usepackage{pgfplots}
\usepgfplotslibrary{groupplots}
\pgfplotsset{compat=1.18}
\usepackage{algorithm}
\usepackage{algpseudocode}
\usepackage[colorlinks=true,linkcolor=blue!50!black,citecolor=blue!50!black,
  urlcolor=blue!50!black]{hyperref}
\pdfstringdefDisableCommands{\def\corref#1{}}  % no footnote mark in the PDF metadata
\usepackage[capitalise]{cleveref}
\crefname{equation}{Eq.}{Eqs.}\Crefname{equation}{Equation}{Equations}
\crefname{section}{Sect.}{Sects.}\Crefname{section}{Section}{Sections}
\crefname{subsection}{Sect.}{Sects.}\Crefname{subsection}{Section}{Sections}
\crefname{figure}{Fig.}{Figs.}\Crefname{figure}{Figure}{Figures}
\crefname{table}{Table}{Tables}\Crefname{table}{Table}{Tables}
\crefname{appendix}{Appendix}{Appendices}\Crefname{appendix}{Appendix}{Appendices}
\crefname{algorithm}{Algorithm}{Algorithms}

% Theorem environments
\newtheorem{proposition}{Proposition}
\newtheorem{lemma}{Lemma}
\newtheorem{corollary}{Corollary}
\theoremstyle{definition}
\newtheorem{definition}{Definition}
\theoremstyle{remark}
\newtheorem{remark}{Remark}

% Notation
\newcommand{\vect}[1]{\bm{#1}}
\newcommand{\matr}[1]{\bm{#1}}
\newcommand{\T}{\top}
\newcommand{\dd}{\,\mathrm{d}}
\newcommand{\bx}{\vect x}
\newcommand{\bxi}{\bm\xi}
\newcommand{\R}{\mathbb{R}}
\newcommand{\Pp}{\mathcal P}
\newcommand{\Qp}{\mathcal Q}
\newcommand{\Tp}{\mathcal T}
\newcommand{\Cp}{\mathcal C}
\newcommand{\Ihat}{\hat\Pi}
\newcommand{\Lam}{\Lambda}
\newcommand{\adj}{\operatorname{adj}}
\newcommand{\spn}{\operatorname{span}}
\newcommand{\cond}{\kappa}
% Notation of the companion paper (Eisentraeger et al., axisymmetric solids)
\newcommand{\Kh}{\hat K}
\newcommand{\bxib}{\bar{\bm\xi}}
\newcommand{\bxit}{\tilde{\bm\xi}}
\newcommand{\Qsk}{\Qp_p(\bar\xi,\bar\eta)}
\newcommand{\nG}{n_{\mathrm{G}}}
\newcommand{\enorm}[1]{\lVert #1\rVert_{\mathrm{E}}}
\newcommand{\enormh}[1]{\lVert #1\rVert_{\mathrm{E},h}}
% Formulations: conventional spectral element, unsymmetric spectral element
% with the metric space Q_p in skew coordinates (USEM of the companion paper),
% unsymmetric spectral element with the hybrid trial space
\newcommand{\PP}{SEM}
\newcommand{\PM}{USEM}
\newcommand{\PH}{USEM-H}
\newcommand{\data}{data}   % path to the data files (relative to main.tex)

% Plot styles
\pgfplotscreateplotcyclelist{bases}{%
  {red,mark=*},{blue,mark=square*},{black,mark=triangle*},
  {green!50!black,mark=diamond*},{orange,mark=pentagon*}}
\pgfplotsset{
  every axis/.append style={tick label style={font=\footnotesize},
    label style={font=\small}, title style={font=\small},
    legend style={font=\footnotesize}},
  styPP/.style={black,mark=o,mark size=1.6pt,thick},
  styPM/.style={red,mark=square*,mark size=1.4pt,thick},
  styPH/.style={blue,mark=triangle*,mark size=1.8pt,thick},
  styLN/.style={black,dashed,thick},
}


% Styles of the thick-cylinder figures (colours for p = 2..8, dark = high p)
\newcommand{\eE}{e_{E}}
\definecolor{p2}{HTML}{FDB863}
\definecolor{p3}{HTML}{F08B3E}
\definecolor{p4}{HTML}{D9541A}
\definecolor{p5}{HTML}{B2182B}
\definecolor{p6}{HTML}{7F1C6A}
\definecolor{p7}{HTML}{4A2A8A}
\definecolor{p8}{HTML}{1B2A5E}
\definecolor{cPP}{HTML}{555555}
\definecolor{cPM}{HTML}{C0392B}
\definecolor{cPH}{HTML}{1F5FA8}
\pgfplotsset{
  conv/.style={
    xmode=log, ymode=log, x dir=reverse, grid=major,
    xtick={1,0.5,0.25,0.125}, xticklabels={$1$,$\frac12$,$\frac14$,$\frac18$},
    xmin=0.1, xmax=1.25, unbounded coords=jump, mark size=1.2pt, line width=0.7pt,
  },
  pline/.style={color=p#1, mark=*, mark options={solid,fill=p#1}},
  meshplot/.style={width=3.0cm, height=3.0cm, axis equal image, hide axis,
    title style={font=\scriptsize, yshift=-1.2ex}},
}

% Avoid overfull lines with long unbreakable math/macros
\setlength{\emergencystretch}{2em}
